\section{Limitations and conclusion}

The measurements support a specific local story.  The large fast fluctuation is the difficulty of the arriving targets under a fixed predictor.  Once that term is removed, the realized finite response is well described on average by the batch's frozen gradient acting on the observed displacement plus the positive curvature cost of moving probabilities.  Linear work generates the loss decrease; softmax curvature offsets part of it.  The optimizer changes both terms, and the network representation residual becomes visible when spans are long or when two optimizer responses are subtracted.

The result does not establish the stronger claim that natural Transformer training universally preserves a constant token-wise optimizer offset.  The scalar probe fails to describe batch-specific response, the optimizer-difference closure fails its registered 10\% RMS gate, and the known-teacher population gap varies strongly across prefixes.  The original RGI phenomenon and the local decomposition can therefore coexist: shared difficulty can produce very high loss correlation while the smaller paired gap remains structured.

The study has four scope limits.  It uses a Pythia-70M local continuation rather than the original paper's 124M and 0.7B from-scratch runs; the three permutations share one data pool; the data's overlap with Pythia pretraining is unknown; and the response diagnostics use observed finite displacements.  The independent checks audit saved arithmetic and fixtures but do not rerun the full gradients or optimizer trajectories.

The practical theoretical target is now concrete.  A future test of strong RGI should measure candidate support mass and within-support shape, record batch-specific gradient response, and prospectively evaluate an unseen window.  It should not infer a universal mechanism from a quadratic sufficient example or from the exact identity \(W+Q\).  Our current evidence identifies the dominant terms in the finite Pythia-70M continuation and an explicit obstruction to extending them to a universal Transformer law.
